% ============================================================
% ГЛАВА: ТЕСТИРОВАНИЕ И ОЦЕНКА
% ============================================================

\section{Тестирование и оценка}

\subsection{Общий подход к тестированию}

Тестирование инструмента выполнялось на двух уровнях: модульном
и интеграционном.
Модульные тесты изолированно проверяют каждый из шести компонентов
системы: загрузку базы данных уязвимых функций, инвентаризацию проекта,
regex"=анализатор, AST"=анализатор, оценщик риска и модуль метрик.
Интеграционные тесты запускают полный конвейер обработки на реальном
тестовом файле и сверяют итоговый план миграции с эталонной разметкой.

Разделение на два уровня обусловлено архитектурой системы: каждый
модуль находится на критическом пути вывода, поэтому ошибка в
отчёте может быть вызвана разными причинами~--- неправильной загрузкой
базы, ошибкой анализатора, ошибкой расчёта риска или ошибкой сериализации
JSON.
Модульные тесты позволяют локализовать источник проблемы, а
интеграционные~--- убедиться, что все части работают совместно и
дают воспроизводимый результат.

Все тесты реализованы на C++ с использованием собственного минималистичного
фреймворка \texttt{TestSuite / TestRunner}, не требующего сторонних
зависимостей.
Итого выполняется 47 тестов: 41 модульный и 6 интеграционных.
Классификация приведена в таблице~\ref{tab:testing_levels}.

\begin{longtable}{|p{4.2cm}|p{4.8cm}|p{6.0cm}|}
  \caption{Классификация тестов системы анализа}
  \label{tab:testing_levels} \\
  \hline
  Уровень & Что проверяется & Примеры тест"=кейсов \\ \hline
  \endfirsthead
  \multicolumn{3}{r}{\textit{Продолжение таблицы \thetable}} \\ \hline
  Уровень & Что проверяется & Примеры тест"=кейсов \\ \hline
  \endhead \hline \endlastfoot

  Модульные тесты базы данных (9)
    & Загрузка JSON, поиск по имени, псевдониму и регистру; слияние с дополнительной базой; наличие ссылок FIPS и ТК\,26
    & \texttt{find\_alias}, \texttt{replacements\_present}, \texttt{merge\_extra\_db} \\ \hline

  Модульные тесты инвентаризации (5)
    & Обход файловой системы, классификация файлов, статистика языков, формирование метаданных CBOM
    & \texttt{cpp\_files\_detected}, \texttt{language\_stats\_cpp}, \texttt{cbom\_metadata\_json} \\ \hline

  Модульные тесты regex"=анализатора (8)
    & Обнаружение функций в тестовом файле, корректность номера строки, отсутствие ложных срабатываний на AES
    & \texttt{finds\_rsa\_in\_fixture}, \texttt{no\_false\_positives\_for\_safe\_code} \\ \hline

  Модульные тесты AST"=анализатора (7)
    & Наличие находок, извлечение контекстной функции и класса, исключение системных заголовков
    & \texttt{context\_function\_extracted}, \texttt{no\_system\_header\_findings} \\ \hline

  Модульные тесты оценщика риска (7)
    & Влияние контекста на итоговую оценку, сортировка плана миграции, наличие рекомендаций и ссылок
    & \texttt{network\_file\_increases\_risk}, \texttt{replacement\_info\_present} \\ \hline

  Модульные тесты метрик (5)
    & Вычисление TP, FP, FN с допуском строки; precision, recall, F1; сериализация в JSON
    & \texttt{perfect\_match}, \texttt{line\_tolerance}, \texttt{json\_serialization} \\ \hline

  Интеграционные тесты (6)
    & Полный конвейер в режимах regex и AST; precision/recall\,$\geq$\,50\,\%; наличие ссылок ТК\,26 и FIPS; сериализация отчёта
    & \texttt{regex\_pipeline\_end\_to\_end}, \texttt{ast\_pipeline\_end\_to\_end}, \texttt{regex\_precision\_recall\_gt50pct} \\

\end{longtable}

% -----------------------------------------------------------
\subsection{Тестовые данные}

Основным тестовым файлом является \texttt{tests/fixtures/sample\_vulnerable.cpp}.
Файл содержит восемь тематических секций (A--H) с заранее известным набором
уязвимых вызовов и отдельный блок \texttt{SAFE}, который не должен давать
срабатываний.
Такой файл выполняет роль контролируемой среды: каждая строка в нём
либо заведомо уязвима, либо заведомо безопасна, что позволяет точно
измерять precision и recall.

Структура тестового файла по секциям представлена в
таблице~\ref{tab:fixture_sections}.

\begin{longtable}{|p{2.2cm}|p{4.5cm}|p{8.3cm}|}
  \caption{Секции тестового файла \texttt{sample\_vulnerable.cpp}}
  \label{tab:fixture_sections} \\
  \hline
  Секция & Класс/контекст & Описание \\ \hline
  \endfirsthead
  \multicolumn{3}{r}{\textit{Продолжение таблицы \thetable}} \\ \hline
  Секция & Класс/контекст & Описание \\ \hline
  \endhead \hline \endlastfoot

  A. \texttt{crypto::tls}   & \texttt{TLSHandshake}    & RSA"=ключ сервера (\texttt{RSA\_new}, \texttt{RSA\_generate\_key\_ex}), ECDH (\texttt{EC\_KEY\_generate\_key}, \texttt{ECDH\_compute\_key}), подпись ECDSA \\ \hline
  B. \texttt{crypto::storage} & \texttt{SecureStorage}  & RSA"=шифрование (\texttt{RSA\_public\_encrypt}), подпись RSA, хеширование SHA"=1 \\ \hline
  C. DH params              & функция \texttt{setup\_dh} & Параметры Диффи"=Хеллмана: \texttt{DH\_new}, \texttt{DH\_generate\_parameters\_ex} \\ \hline
  D. MD5                    & функция \texttt{md5sum}  & Устаревшая хеш"=функция MD5 \\ \hline
  E. \texttt{crypto::evp}   & \texttt{EvpOperations}   & EVP высокого уровня: контекст с классическим алгоритмом, подпись и верификация через EVP, ключевое согласование ECDH через EVP \\ \hline
  F. \texttt{crypto::dsa\_legacy} & \texttt{DsaOps}    & Явные вызовы DSA (\texttt{DSA\_new}, \texttt{DSA\_sign}); проверяет, что анализатор не смешивает DSA и ECDSA \\ \hline
  G. \texttt{crypto::ec\_ops} & \texttt{ECOperations}  & Низкоуровневые EC"=операции: создание группы, \texttt{EC\_KEY\_new}, \texttt{EC\_KEY\_generate\_key} \\ \hline
  H. \texttt{crypto::bn\_ops} & \texttt{BNOperations}  & Ручные BN"=операции как индикатор ручной реализации RSA: \texttt{BN\_generate\_prime\_ex}, \texttt{BN\_mod\_exp} \\ \hline
  SAFE (нет срабатываний)   & функции \texttt{safe\_*} & AES"=256"=GCM (симметричный, допустим); \texttt{EVP\_PKEY\_CTX\_new\_from\_name} с аргументом \texttt{"ML-KEM-768"} (постквантовый) \\

\end{longtable}

Полный набор ожидаемых находок из 32 записей хранится в файле
\texttt{ground\_truth.json}.
Каждая запись содержит имя функции, номер строки и допуск
(\texttt{line\_tolerance}~--- обычно 3--5 строк), позволяющий
корректно засчитывать совпадение при незначительных расхождениях
номеров строк между AST и regex"=режимами.
Перечень ожидаемых функций, сгруппированных по категориям, приведён в
таблице~\ref{tab:expected_functions}.

\begin{longtable}{|p{1.2cm}|p{5.8cm}|p{7.9cm}|}
  \caption{Ожидаемые находки в тестовом файле (ground truth, 32 записи)}
  \label{tab:expected_functions} \\
  \hline
  Стр. & Функция & Обоснование обнаружения \\ \hline
  \endfirsthead
  \multicolumn{3}{r}{\textit{Продолжение таблицы \thetable}} \\ \hline
  Стр. & Функция & Обоснование обнаружения \\ \hline
  \endhead \hline \endlastfoot

  49  & \texttt{RSA\_new}                   & Создание RSA-структуры; индикатор использования квантово-уязвимого RSA \\ \hline
  52  & \texttt{RSA\_generate\_key\_ex}      & Генерация RSA-ключа; алгоритм полностью сломан алгоритмом Шора \\ \hline
  60  & \texttt{EC\_KEY\_generate\_key}      & Генерация EC-ключа; ECDLP решается за полиномиальное время \\ \hline
  61  & \texttt{ECDH\_compute\_key}          & Согласование ключа на эллиптических кривых \\ \hline
  69  & \texttt{ECDSA\_sign}                 & Подпись ECDSA; основана на ECDLP \\ \hline
  83  & \texttt{RSA\_public\_encrypt}        & RSA-шифрование; уязвимо к атаке Шора \\ \hline
  89  & \texttt{RSA\_sign}                   & RSA-подпись; уязвима к атаке Шора \\ \hline
  94  & \texttt{SHA1}                        & SHA-1: коллизии найдены классически; Гровер снижает стойкость до~\(\approx80\)\,бит \\ \hline
  103 & \texttt{DH\_new}                     & Создание DH-параметров; задача дискретного логарифма решается алгоритмом Шора \\ \hline
  104 & \texttt{DH\_generate\_parameters\_ex}& Генерация DH-параметров; псевдоним в базе для группы \texttt{DH\_generate\_key} \\ \hline
  110 & \texttt{MD5}                         & MD5: коллизии полностью сломаны; Гровер даёт~\(\approx64\)\,бит прообраза \\ \hline
  121 & \texttt{EVP\_PKEY\_CTX\_new\_from\_name} & Аргумент \texttt{"RSA"} --- классический алгоритм; условно уязвим \\ \hline
  124 & \texttt{EVP\_PKEY\_keygen\_init}     & Контекст инициализирован для RSA-ключа; уязвим \\ \hline
  125 & \texttt{EVP\_PKEY\_keygen}           & Генерация ключа в RSA-контексте \\ \hline
  143 & \texttt{EVP\_PKEY\_CTX\_new}         & Создание контекста из существующего EC/RSA-ключа; безусловно уязвим \\ \hline
  150 & \texttt{EVP\_DigestSignInit}         & Инициализация операции подписи; уязвима при RSA/ECDSA-ключе \\ \hline
  152 & \texttt{EVP\_DigestSignUpdate}       & Передача данных в операцию подписи \\ \hline
  154 & \texttt{EVP\_DigestSignFinal}        & Завершение подписи; финальная операция \\ \hline
  162 & \texttt{EVP\_DigestVerifyInit}       & Инициализация верификации; уязвима при RSA/ECDSA-ключе \\ \hline
  164 & \texttt{EVP\_DigestVerifyUpdate}     & Передача данных в операцию верификации \\ \hline
  165 & \texttt{EVP\_DigestVerifyFinal}      & Завершение верификации \\ \hline
  171 & \texttt{EVP\_PKEY\_CTX\_new}         & Второй экземпляр: ECDH-контекст через существующий ключ \\ \hline
  172 & \texttt{EVP\_PKEY\_derive\_init}     & Инициализация ECDH-согласования через EVP \\ \hline
  173 & \texttt{EVP\_PKEY\_derive\_set\_peer}& Установка открытого ключа партнёра \\ \hline
  175 & \texttt{EVP\_PKEY\_derive}           & Выработка общего секрета ECDH \\ \hline
  186 & \texttt{DSA\_new}                    & Создание DSA-структуры; DSA основан на задаче дискретного логарифма \\ \hline
  191 & \texttt{DSA\_sign}                   & DSA-подпись; уязвима к алгоритму Шора \\ \hline
  201 & \texttt{EC\_GROUP\_new\_by\_curve\_name} & Выбор эллиптической кривой; все именованные кривые NIST квантово-уязвимы \\ \hline
  203 & \texttt{EC\_KEY\_new}                & Выделение EC-структуры \\ \hline
  205 & \texttt{EC\_KEY\_generate\_key}      & Второй экземпляр генерации EC-ключа в той же функции \\ \hline
  218 & \texttt{BN\_generate\_prime\_ex}     & Генерация большого простого числа; индикатор ручной реализации RSA \\ \hline
  225 & \texttt{BN\_mod\_exp}                & Модульное возведение в степень; ключевая операция RSA/DH \\ \hline

\end{longtable}

Блок \texttt{SAFE} содержит два фрагмента, которые не должны приводить
к срабатыванию.
Первый~--- функция \texttt{safe\_aes\_encrypt}, использующая
\texttt{EVP\_EncryptInit\_ex} с \texttt{EVP\_aes\_256\_gcm()}: симметричный
режим AES-256-GCM не является объектом поиска квантово-уязвимых
асимметричных алгоритмов и не входит в базу.
Второй~--- функция \texttt{safe\_pqc\_kem}, вызывающая
\texttt{EVP\_PKEY\_CTX\_new\_from\_name} с аргументом
\texttt{"ML-KEM-768"}: AST-анализатор прослеживает провенанс
аргумента и подавляет находку при обнаружении постквантового имени
алгоритма, тогда как regex-анализатор в данном случае даёт ложное
срабатывание (подробнее~--- в разделе~\ref{sec:analysis_comparison}).

% -----------------------------------------------------------
\subsection{Модульные тесты}

\subsubsection{Тесты базы данных уязвимых функций}

Набор из девяти тестов класса \texttt{VulnDatabase} проверяет загрузку
JSON-файла, поиск по основному имени, псевдониму, нечувствительности
к регистру, отсутствие ответа для несуществующей функции, наличие
заполненных рекомендаций по миграции и ссылок на стандарты, а
также слияние с дополнительной базой.

Тест поиска по каноническому имени (листинг~\ref{lst:db-find-rsa})
проверяет, что запись для \texttt{RSA\_generate\_key\_ex} загружена,
алгоритм равен \texttt{"RSA"}, уязвимость~--- \texttt{"high"}, а оценка
риска превышает 9.0.

\begin{itmolisting}[C++]{Тест поиска RSA-функции в базе данных}{lst:db-find-rsa}
TEST(VulnDatabase, find_rsa_keygen) {
    pqc::VulnDatabase db;
    db.load("data/vulnerable_functions.json");

    auto opt = db.find_by_name("RSA_generate_key_ex");

    ASSERT_TRUE(opt.has_value());
    ASSERT_EQ(opt->algorithm, "RSA");
    ASSERT_EQ(opt->quantum_vulnerability, "high");
    ASSERT_GT(opt->risk_score, 9.0);
}
\end{itmolisting}

Тест псевдонима (листинг~\ref{lst:db-find-alias}) важен, поскольку
в базе данных у одной записи может быть несколько совместимых имён~---
например, \texttt{RSA\_generate\_key} является псевдонимом канонической
функции \texttt{RSA\_generate\_key\_ex}.
Тест проверяет, что поиск по псевдониму возвращает ту же запись.

\begin{itmolisting}[C++]{Тест поиска по псевдониму функции}{lst:db-find-alias}
TEST(VulnDatabase, find_alias) {
    pqc::VulnDatabase db;
    db.load("data/vulnerable_functions.json");

    auto opt = db.find_by_name("RSA_generate_key"); // псевдоним

    ASSERT_TRUE(opt.has_value());
    ASSERT_EQ(opt->name, "RSA_generate_key_ex");    // каноническое имя
}
\end{itmolisting}

Тест слияния с дополнительной базой (листинг~\ref{lst:db-merge})
моделирует сценарий, когда организация добавляет собственные
функции"=обёртки для корпоративных криптографических API.

\begin{itmolisting}[C++]{Тест слияния с дополнительной базой}{lst:db-merge}
TEST(VulnDatabase, merge_extra_db) {
    pqc::VulnDatabase db;
    db.load("data/vulnerable_functions.json");
    size_t before = db.size();

    nlohmann::json extra = {
      {"version", "1.0"},
      {"libraries", {{"internal", {{"functions", nlohmann::json::array({
        {{"id","corp-1"},{"name","CorpRSAEncrypt"},
         {"aliases", nlohmann::json::array()},
         {"category","asymmetric_encryption"},
         {"algorithm","RSA"},{"quantum_vulnerability","high"},
         {"risk_score",9.5},{"description","Internal RSA wrapper"},
         {"patterns", nlohmann::json::array({"CorpRSAEncrypt\\s*\\("})},
         {"replacements", nlohmann::json::array()}}
      })}}}}}
    };

    db.merge_json(extra);

    ASSERT_GT((int)db.size(), (int)before);
    ASSERT_TRUE(db.find_by_name("CorpRSAEncrypt").has_value());
}
\end{itmolisting}

Тест \texttt{replacements\_present} дополнительно проверяет, что первая
рекомендация по миграции для \texttt{RSA\_generate\_key\_ex} содержит
ссылку на стандарт FIPS (поле \texttt{standard}), что обеспечивает
корректную сериализацию плана миграции в JSON-отчёте.

\subsubsection{Тесты инвентаризации проекта}

Набор из пяти тестов класса \texttt{ProjectScanner} проверяет, что сканер
обходит директорию, находит C++"=файлы, корректно определяет язык,
формирует метаданные CBOM и поддерживает пользовательские плагины
языка (листинг~\ref{lst:scanner-cpp}).

\begin{itmolisting}[C++]{Тест обнаружения C++ файлов при инвентаризации}{lst:scanner-cpp}
TEST(ProjectScanner, cpp_files_detected) {
    pqc::ProjectScanner scanner;
    auto inv = scanner.scan("tests/fixtures");

    auto src = inv.get_source_files();

    ASSERT_GT((int)src.size(), 0);
    bool found_cpp = false;
    for (auto* f : src)
        if (f->path.extension() == ".cpp") { found_cpp = true; break; }
    ASSERT_TRUE(found_cpp);
}
\end{itmolisting}

Тест \texttt{custom\_plugin} демонстрирует расширяемость: в сканер
регистрируется экземпляр анонимного класса, реализующего интерфейс
\texttt{ILanguagePlugin}, и проверяется, что инвентаризация проходит
без ошибок.

\subsubsection{Тесты regex"=анализатора}

Восемь тестов класса \texttt{RegexAnalyzer} проверяют обнаружение
конкретных функций, атрибуты находки (номер строки, библиотека, ссылка
NIST) и отсутствие ложных срабатываний на AES.
Ключевые тест"=кейсы представлены ниже.

Тест \texttt{finds\_rsa\_in\_fixture} (листинг~\ref{lst:regex-rsa})
проверяет, что анализатор обнаруживает вызов \texttt{RSA\_generate\_key\_ex}
в тестовом файле.

\begin{itmolisting}[C++]{Тест обнаружения RSA-функции regex-анализатором}{lst:regex-rsa}
TEST(RegexAnalyzer, finds_rsa_in_fixture) {
    pqc::VulnDatabase db;
    db.load("data/vulnerable_functions.json");
    pqc::RegexAnalyzer ra(db);

    auto findings = ra.analyze_file(
        "tests/fixtures/sample_vulnerable.cpp"
    );

    bool found = false;
    for (auto& f : findings)
        if (f.function_name == "RSA_generate_key_ex")
            { found = true; break; }
    ASSERT_TRUE(found);
}
\end{itmolisting}

Тест \texttt{no\_false\_positives\_for\_safe\_code}
(листинг~\ref{lst:regex-safe}) проверяет, что симметричный вызов
\texttt{EVP\_EncryptInit\_ex} в блоке \texttt{SAFE} не порождает
находок~--- функция отсутствует в базе уязвимых функций.

\begin{itmolisting}[C++]{Тест отсутствия ложных срабатываний на AES-код}{lst:regex-safe}
TEST(RegexAnalyzer, no_false_positives_for_safe_code) {
    pqc::VulnDatabase db;
    db.load("data/vulnerable_functions.json");
    pqc::RegexAnalyzer ra(db);
    auto findings = ra.analyze_file(
        "tests/fixtures/sample_vulnerable.cpp"
    );

    for (auto& f : findings)
        ASSERT_TRUE(f.function_name != "EVP_EncryptInit_ex");
}
\end{itmolisting}

Тест \texttt{finding\_has\_nist\_reference} перебирает все находки и
проверяет, что каждая из них содержит непустую строку ссылки на стандарт
NIST~--- это обязательный атрибут для формирования полноценного
раздела \texttt{nist\_compliance} в отчёте.

\subsubsection{Тесты AST"=анализатора}

Семь тестов класса \texttt{ASTAnalyzer} проверяют режим работы, наличие
libclang, обнаружение уязвимостей и корректность извлечения контекста.

Тест \texttt{context\_function\_extracted}
(листинг~\ref{lst:ast-ctx-fn}) проверяет, что хотя бы для одной находки
поле \texttt{context\_function} непусто и не равно \texttt{"<global>"}~---
то есть AST-анализатор действительно определяет объемлющую функцию.

\begin{itmolisting}[C++]{Тест извлечения контекстной функции AST-анализатором}{lst:ast-ctx-fn}
TEST(ASTAnalyzer, context_function_extracted) {
    pqc::VulnDatabase db;
    db.load("data/vulnerable_functions.json");
    pqc::ASTAnalyzer aa(db);
    auto findings = aa.analyze_file(
        "tests/fixtures/sample_vulnerable.cpp"
    );

    bool has_context = false;
    for (auto& f : findings)
        if (!f.context_function.empty() &&
             f.context_function != "<global>")
            { has_context = true; break; }
    ASSERT_TRUE(has_context);
}
\end{itmolisting}

Тест \texttt{context\_class\_extracted} аналогично проверяет, что для
хотя бы одной находки заполнено поле \texttt{context\_class}, что
подтверждает способность анализатора определять вмещающий класс (например,
\texttt{crypto::tls::TLSHandshake}).

Тест \texttt{no\_system\_header\_findings} (листинг~\ref{lst:ast-nosys})
критичен для корректности: он проверяет, что все находки указывают
на файл \texttt{sample\_vulnerable.cpp}, а не на системные заголовки
OpenSSL, подключаемые при разборе.

\begin{itmolisting}[C++]{Тест исключения находок из системных заголовков}{lst:ast-nosys}
TEST(ASTAnalyzer, no_system_header_findings) {
    pqc::VulnDatabase db;
    db.load("data/vulnerable_functions.json");
    pqc::ASTAnalyzer aa(db);
    auto findings = aa.analyze_file(
        "tests/fixtures/sample_vulnerable.cpp"
    );

    for (auto& f : findings)
        ASSERT_CONTAINS(f.file_path, "sample_vulnerable");
}
\end{itmolisting}

\subsubsection{Тесты оценщика риска}

Семь тестов класса \texttt{RiskAssessor} проверяют, что контекстные
факторы корректно влияют на итоговую оценку.

Тест \texttt{test\_file\_reduces\_risk}
(листинг~\ref{lst:risk-test-reduce}) подтверждает, что путь
\texttt{tests/test\_tls.cpp} распознаётся как тестовый файл и
применяется понижающий коэффициент.

\begin{itmolisting}[C++]{Тест снижения риска для тестовых файлов}{lst:risk-test-reduce}
TEST(RiskAssessor, test_file_reduces_risk) {
    pqc::VulnDatabase db;
    db.load("data/vulnerable_functions.json");
    pqc::RiskAssessor ra(db);
    pqc::ProjectInventory inv;

    pqc::Finding f;
    f.function_name  = "RSA_generate_key_ex";
    f.library        = "openssl";
    f.category       = "asymmetric_key_generation";
    f.base_risk_score = 9.5;
    f.file_path       = "tests/test_tls.cpp";

    auto rr = ra.assess({f}, inv);

    ASSERT_NOT_EMPTY(rr.migration_plan);
    // Тестовый файл применяет коэффициент x0.7
    ASSERT_TRUE(rr.migration_plan[0].risk.final_score < 9.5);
}
\end{itmolisting}

Тест \texttt{migration\_order\_sorted\_by\_risk}
(листинг~\ref{lst:risk-order}) намеренно подаёт две находки в
обратном порядке и проверяет, что план миграции отсортирован по
убыванию итоговой оценки риска.

\begin{itmolisting}[C++]{Тест сортировки плана миграции по риску}{lst:risk-order}
TEST(RiskAssessor, migration_order_sorted_by_risk) {
    pqc::VulnDatabase db;
    db.load("data/vulnerable_functions.json");
    pqc::RiskAssessor ra(db);
    pqc::ProjectInventory inv;

    pqc::Finding f1, f2;
    // f1 высокий риск (9.5), f2 средний (6.5)
    f1.function_name = "RSA_generate_key_ex"; f1.base_risk_score = 9.5;
    f1.file_path = "src/net.cpp";  f1.category = "asymmetric_key_generation";
    f2.function_name = "SHA1";     f2.base_risk_score = 6.5;
    f2.file_path = "src/util.cpp"; f2.category = "hash_function";

    // Подаём намеренно в обратном порядке
    auto rr = ra.assess({f2, f1}, inv);

    ASSERT_EQ((int)rr.migration_plan.size(), 2);
    ASSERT_GE(rr.migration_plan[0].risk.final_score,
              rr.migration_plan[1].risk.final_score);
}
\end{itmolisting}

\subsubsection{Тесты модуля метрик}

Пять тестов класса \texttt{MetricsCalculator} проверяют вычисление TP,
FP, FN, precision, recall и F1, а также работу допуска строк и
сериализацию результатов в JSON.

Тест \texttt{line\_tolerance} (листинг~\ref{lst:metrics-tol}) проверяет,
что находка на строке~45 засчитывается как TP при эталонном значении
строки~42 и допуске~5.

\begin{itmolisting}[C++]{Тест допуска при сопоставлении строк}{lst:metrics-tol}
TEST(MetricsCalculator, line_tolerance) {
    pqc::MetricsCalculator calc;

    pqc::Finding f;
    f.function_name = "RSA_sign";
    f.line_number   = 45;

    pqc::GroundTruth gt;
    gt.function_name = "RSA_sign";
    gt.line_number   = 42;
    gt.line_tolerance = 5;

    auto m = calc.compute({f}, {gt});

    ASSERT_EQ(m.true_positives, 1); // в пределах допуска
    ASSERT_EQ(m.false_positives, 0);
    ASSERT_EQ(m.false_negatives, 0);
}
\end{itmolisting}

% -----------------------------------------------------------
\subsection{Интеграционные тесты}

Шесть интеграционных тестов запускают полный конвейер обработки на
директории \texttt{tests/fixtures} и проверяют корректность итогового
результата.

Тест \texttt{regex\_pipeline\_end\_to\_end}
(листинг~\ref{lst:itest-regex}) последовательно выполняет все
этапы: загрузку базы, инвентаризацию, regex"=анализ и оценку риска,
после чего проверяет непустоту плана миграции и ненулевую общую оценку
риска.

\begin{itmolisting}[C++]{Интеграционный тест полного конвейера в режиме regex}{lst:itest-regex}
TEST(FullPipeline, regex_pipeline_end_to_end) {
    pqc::VulnDatabase db;
    db.load("data/vulnerable_functions.json");

    pqc::ProjectScanner scanner;
    auto inv = scanner.scan("tests/fixtures");

    pqc::RegexAnalyzer ra(db);
    auto findings = ra.analyze(inv);
    ASSERT_NOT_EMPTY(findings);

    pqc::RiskAssessor assessor(db);
    auto rr = assessor.assess(findings, inv);

    ASSERT_NOT_EMPTY(rr.migration_plan);
    ASSERT_GT(rr.overall_risk_score, 0.0);
    ASSERT_NOT_EMPTY(rr.nist_readiness);
}
\end{itmolisting}

Тест \texttt{ast\_pipeline\_end\_to\_end} аналогично проверяет AST"=режим
и выводит в лог число найденных уязвимостей и используемый бэкенд
(libclang или token"=fallback).

Тест \texttt{regex\_precision\_recall\_gt50pct}
(листинг~\ref{lst:itest-metrics}) является регрессионным: он загружает
эталонную разметку из \texttt{ground\_truth.json} и проверяет, что
precision, recall и F1 не опускаются ниже~50\,\%.
Порог выбран намеренно консервативным, чтобы тест не зависел от
конкретного состояния базы, однако фиксировал регресс при значительном
ухудшении.

\begin{itmolisting}[C++]{Интеграционный тест метрик качества}{lst:itest-metrics}
TEST(FullPipeline, regex_precision_recall_gt50pct) {
    pqc::VulnDatabase db;
    db.load("data/vulnerable_functions.json");

    pqc::ProjectScanner scanner;
    auto inv = scanner.scan("tests/fixtures");

    pqc::RegexAnalyzer ra(db);
    auto findings = ra.analyze(inv);

    auto gt = pqc::MetricsCalculator::load_from_json(
        "tests/fixtures/ground_truth.json"
    );
    pqc::MetricsCalculator calc;
    auto m = calc.compute(findings, gt);

    ASSERT_GE(m.precision(), 0.5);
    ASSERT_GE(m.recall(),    0.5);
    ASSERT_GE(m.f1_score(),  0.5);
}
\end{itmolisting}

Тест \texttt{migration\_plan\_has\_fips\_references} проверяет, что
хотя бы в одной записи плана миграции поле \texttt{replacement\_standard}
содержит подстроку \texttt{"FIPS"}, что подтверждает корректную
загрузку и применение рекомендаций из базы.

% -----------------------------------------------------------
\subsection{Оценка результатов тестирования}

Все~47 тестов~--- 41 модульный и 6 интеграционных~--- проходят
успешно (таблица~\ref{tab:test_results_summary}).
До внесения исправлений, описанных в разделе~\ref{sec:analysis_comparison},
три теста завершались с ошибкой:
\texttt{replacements\_present} (поле \texttt{standard} было пустым
из"=за расхождения имён полей в JSON),
\texttt{replacement\_info\_present} и \texttt{tc26\_note\_present}
(поля заполнялись только из объекта \texttt{Finding}, но тестовые
находки, созданные вручную, не несли этих атрибутов).
Исправления состояли в расширении парсера базы данных для принятия
обоих форматов полей (\texttt{"function"}/\texttt{"function\_name"},
\texttt{"notes"}/\texttt{"migration\_notes"}) и добавлении запасного
обращения к базе при пустых полях в объекте \texttt{MigrationAction}.

\begin{longtable}{|p{5.5cm}|r|r|r|}
  \caption{Итоги выполнения тестов}
  \label{tab:test_results_summary} \\
  \hline
  Набор тестов & Всего & Пройдено & Провалено \\ \hline
  \endfirsthead
  \multicolumn{4}{r}{\textit{Продолжение таблицы \thetable}} \\ \hline
  Набор тестов & Всего & Пройдено & Провалено \\ \hline
  \endhead \hline \endlastfoot

  VulnDatabase      & 9  & 9  & 0 \\ \hline
  ProjectScanner    & 5  & 5  & 0 \\ \hline
  RegexAnalyzer     & 8  & 8  & 0 \\ \hline
  ASTAnalyzer       & 7  & 7  & 0 \\ \hline
  RiskAssessor      & 7  & 7  & 0 \\ \hline
  MetricsCalculator & 5  & 5  & 0 \\ \hline
  FullPipeline (интеграционные) & 6 & 6 & 0 \\ \hline
  \textbf{Итого}    & \textbf{47} & \textbf{47} & \textbf{0} \\

\end{longtable}

Интеграционный тест \texttt{regex\_precision\_recall\_gt50pct} зафиксировал
финальные метрики качества на эталонном наборе: precision\,=\,80,0\,\%,
recall\,=\,100,0\,\%, F1\,=\,88,9\,\% для режима regex.
AST"=режим даёт precision\,=\,94,1\,\%, recall\,=\,100,0\,\%,
F1\,=\,97,0\,\%.
Подробный сравнительный анализ приведён в разделе~\ref{sec:analysis_comparison}.


% ============================================================
% ГЛАВА: АНАЛИЗ ПОЛУЧЕННЫХ РЕЗУЛЬТАТОВ
% ============================================================

\section{Анализ полученных результатов}

\subsection{Соответствие реализации поставленным требованиям}

В таблице~\ref{tab:req_compliance} приведено соответствие реализованных
функций ключевым требованиям, сформулированным в техническом задании.

\begin{longtable}{|p{6.0cm}|p{4.5cm}|p{4.5cm}|}
  \caption{Соответствие реализации функциональным требованиям}
  \label{tab:req_compliance} \\
  \hline
  Требование & Реализация & Подтверждение \\ \hline
  \endfirsthead
  \multicolumn{3}{r}{\textit{Продолжение таблицы \thetable}} \\ \hline
  Требование & Реализация & Подтверждение \\ \hline
  \endhead \hline \endlastfoot

  Обнаружение квантово-уязвимых функций в C/C++ коде
    & Два анализатора: regex и AST (libclang)
    & TP\,=\,32/32 на эталонном наборе в обоих режимах \\ \hline

  Поддержка баз OpenSSL, Crypto++, ГОСТ"=криптографии
    & Три библиотечных секции в \texttt{vulnerable\_functions.json}
    & Тест \texttt{gost\_library\_entries} \\ \hline

  Расчёт риска с учётом контекста
    & Модуль \texttt{RiskAssessor}: 6 контекстных множителей
    & Тесты \texttt{network\_file\_increases\_risk}, \texttt{test\_file\_reduces\_risk} \\ \hline

  Формирование плана миграции, отсортированного по приоритету
    & \texttt{RiskAssessor::assess} с сортировкой по \texttt{final\_score}
    & Тест \texttt{migration\_order\_sorted\_by\_risk} \\ \hline

  Рекомендации по миграции со ссылками FIPS и ТК\,26
    & Поля \texttt{replacement\_standard}, \texttt{tc26\_note} в \texttt{MigrationAction}
    & Тесты \texttt{migration\_plan\_has\_fips\_references}, \texttt{risk\_report\_has\_tc26\_data} \\ \hline

  Формирование CBOM по стандарту CycloneDX
    & Модуль \texttt{Cbom} и \texttt{ReportGenerator::build\_cbom}
    & Тест \texttt{report\_json\_serialization} \\ \hline

  Расширяемость через пользовательскую базу
    & Метод \texttt{VulnDatabase::merge\_json}
    & Тест \texttt{merge\_extra\_db} \\ \hline

  Измерение точности и полноты с эталонной разметкой
    & Модуль \texttt{MetricsCalculator}, флаг \texttt{{-}{-}metrics}
    & Тест \texttt{regex\_precision\_recall\_gt50pct} \\

\end{longtable}

Все восемь функциональных требований реализованы и подтверждены
автоматическими тестами.

% -----------------------------------------------------------
\subsection{Сравнение режимов анализа}
\label{sec:analysis_comparison}

Оба режима проверялись на одном и том же тестовом файле
\texttt{sample\_vulnerable.cpp} (32 записи в эталонной разметке).
Результаты представлены в таблице~\ref{tab:mode_comparison}.

\begin{longtable}{|p{4.5cm}|r|r|}
  \caption{Сравнение метрик качества AST и regex"=режимов}
  \label{tab:mode_comparison} \\
  \hline
  Метрика & AST & Regex \\ \hline
  \endfirsthead
  \multicolumn{3}{r}{\textit{Продолжение таблицы \thetable}} \\ \hline
  Метрика & AST & Regex \\ \hline
  \endhead \hline \endlastfoot

  Эталонных записей (GT)    & 32 & 32 \\ \hline
  True Positives (TP)        & 32 & 32 \\ \hline
  False Positives (FP)       & 2  & 8  \\ \hline
  False Negatives (FN)       & 0  & 0  \\ \hline
  Precision                  & 94,1\,\% & 80,0\,\% \\ \hline
  Recall                     & 100,0\,\% & 100,0\,\% \\ \hline
  F1                         & 97,0\,\% & 88,9\,\% \\ \hline
  Время анализа тестового файла & $\approx$1030\,мс & $\approx$80\,мс \\

\end{longtable}

\paragraph{True Positives.}
Оба режима обнаружили все~32 ожидаемых вызова.
В числе найденных~--- условно уязвимый
\texttt{EVP\_PKEY\_CTX\_new\_from\_name("RSA", ...)}, шесть функций
EVP"=подписи/верификации, хеш"=функции SHA"=1 и MD5, низкоуровневые
BIGNUM"=операции, а также псевдонимные вызовы
\texttt{EC\_KEY\_generate\_key} и \texttt{DH\_generate\_parameters\_ex},
которые не были в базе изначально.

\paragraph{False Positives.}
\textbf{AST (2 FP):} обе находки~--- \texttt{BN\_new} на строках~50 и~217~---
являются реальными уязвимостями (создание BIGNUM"=структуры как индикатор
ручной реализации RSA/DH), которые просто отсутствуют в эталонной разметке.
С формальной точки зрения они считаются FP, однако с содержательной~---
это корректные срабатывания.
Если включить эти строки в GT, precision AST достигнет 100\,\%.

\textbf{Regex (8 FP):} включают те же~2 валидных \texttt{BN\_new} плюс~6
ложных срабатываний, обусловленных фундаментальным ограничением подхода:
\begin{itemize}
  \item 3 вызова из блока \texttt{safe\_pqc\_kem}~---
        \texttt{EVP\_PKEY\_CTX\_new\_from\_name("ML-KEM-768", ...)},
        \texttt{EVP\_PKEY\_keygen\_init}, \texttt{EVP\_PKEY\_keygen}~---
        являются постквантовыми и не должны давать находок, однако
        regex"=анализатор не может отследить провенанс аргумента
        \texttt{"ML-KEM-768"};
  \item аналогичные~3 вызова из второго экземпляра \texttt{EVP\_PKEY\_CTX\_new\_from\_name}
        с тем же аргументом.
\end{itemize}
AST"=анализатор корректно подавляет все эти~6 находок, прослеживая цепочку
значений через механизм \texttt{resolve\_expr} и проверяя, содержит ли
аргумент имя постквантового алгоритма.

\paragraph{False Negatives.}
Ни одного FN в обоих режимах.
Нулевое значение FN достигнуто после внесения следующих исправлений в
базу данных:
добавление записей \texttt{SHA1}, \texttt{MD5};
добавление псевдонима \texttt{EC\_KEY\_generate\_key} к \texttt{EC\_KEY\_new};
добавление псевдонима \texttt{DH\_generate\_parameters\_ex} к \texttt{DH\_generate\_key};
смена канонического имени группы RSA-генерации с \texttt{RSA\_generate\_key}
на \texttt{RSA\_generate\_key\_ex}.
Помимо этого, в AST"=анализаторе исправлено поле \texttt{function\_name}:
теперь в находке указывается фактически вызванная функция, а не
каноническое имя в базе, что обеспечивает корректное совпадение с GT.

\paragraph{Рекомендации по применению режимов.}
Regex"=режим целесообразно использовать как быстрый скрининг: время
работы на тестовом файле составляет~$\approx$80\,мс против
$\approx$1030\,мс у AST"=режима.
При этом~6 из~8 FP неустранимы без анализа аргументов.
AST"=режим рекомендуется как основной для итогового отчёта и
для случаев, когда необходимо подавить ложные срабатывания на
постквантовый код.

% -----------------------------------------------------------
\subsection{Оценка практической применимости}

Инструмент был запущен на нескольких открытых проектах с использованием
OpenSSL для оценки практической применимости.
В проектах обнаруживались те же паттерны, что и в тестовом файле:
прямые вызовы \texttt{RSA\_generate\_key\_ex}, \texttt{ECDH\_compute\_key},
блоки EVP"=инициализации.
Время анализа файла объёмом~$\approx$2000\,строк составило
менее~2\,с в AST"=режиме, что является приемлемым для CI"=конвейера.

Помимо поиска уязвимостей, инструмент формирует:
\begin{itemize}
  \item CBOM"=файл в формате CycloneDX с полями \texttt{primitive},
        \texttt{cryptoFunctions}, \texttt{quantumSecurityLevel}~---
        стандартизованный инвентарь криптографических активов;
  \item план миграции с явными ссылками на FIPS\,203/204/205 и
        рекомендации ТК\,26 в переходный период 2025--2030;
  \item оценку NIST"=готовности (\texttt{not-started} / \texttt{in-progress} /
        \texttt{near-compliant} / \texttt{compliant}).
\end{itemize}
Эти выходные данные пригодны как для внутреннего аудита, так и для
представления регулятору.

% -----------------------------------------------------------
\subsection{Ограничения реализации}

\paragraph{Межпроцедурный анализ потоков данных.}
Если ключевой материал передаётся через несколько функций, AST"=анализатор
отслеживает цепочку только в рамках одной функции (через
\texttt{resolve\_expr}). Межпроцедурный анализ не реализован, что может
приводить к пропуску уязвимостей при глубоком стеке вызовов.

\paragraph{Провенанс аргументов в режиме regex.}
Regex"=анализатор не может отслеживать значения переменных, поэтому
не способен подавить ложные срабатывания на постквантовый код
(6~FP из~8 в тестовом наборе). Частичное решение~--- доэмпирические
правила~--- хрупко и не масштабируется.

\paragraph{Зависимость от libclang.}
AST"=анализатор требует установленной libclang и правильно настроенных
путей включения. При отсутствии libclang система переключается на
токенизирующий fallback"=режим с пониженной точностью извлечения контекста.

\paragraph{Покрытие базы данных.}
Текущая база охватывает три библиотеки: OpenSSL, Crypto++ и библиотеку
ГОСТ"=криптографии. Функции из wolfSSL, Botan, mbedTLS и прочих
библиотек не включены и будут пропущены.

\paragraph{Отсутствие модуля автоматической миграции.}
Инструмент формирует план миграции, но не выполняет кодовые
преобразования автоматически. Разработчик должен вручную применять
рекомендации на основе сгенерированного отчёта.

% -----------------------------------------------------------
\subsection{Направления дальнейшего развития}

\paragraph{Модуль автоматической миграции кода.}
Приоритетным направлением является реализация модуля, который по
готовому плану миграции и шаблонам из базы данных генерирует
патчи или заготовки замены~--- например, заменяет
\texttt{RSA\_generate\_key\_ex} на блок инициализации
\texttt{EVP\_PKEY\_CTX\_new\_from\_name("ML-DSA-65", ...)} с
сохранением семантики.
Для этого потребуется расширить AST"=обходчик до уровня переписывающего
трансформера, использующего тот же cursor"=механизм libclang.

\paragraph{Межпроцедурный анализ.}
Реализация простого межпроцедурного анализа потоков данных (хотя бы
в рамках одной единицы трансляции) позволит правильно
оценивать риск при передаче ключевого материала через вспомогательные
функции, что снизит число FN на реальных проектах с абстракциями
над криптографическими API.

\paragraph{Расширение базы данных.}
Добавление функций из wolfSSL, mbedTLS, Botan, а также Python"=биндингов
(\texttt{cryptography}, \texttt{pyOpenSSL}), Go"=стандартной библиотеки
(\texttt{crypto/rsa}, \texttt{crypto/elliptic}) и Rust"=крейтов
(\texttt{rsa}, \texttt{p256}) повысит применимость инструмента для
многоязычных проектов.

\paragraph{Интеграция в CI/CD.}
Добавление режима выхода с ненулевым кодом при обнаружении критических
находок (\texttt{{-}{-}fail-on-critical}) позволит встраивать инструмент
в GitHub Actions, GitLab CI и аналогичные системы как шаг проверки
перед слиянием.

\paragraph{Отслеживание истечения срока действия сертификатов.}
Текущий модуль анализа сертификатов определяет алгоритм, но не
сигнализирует о скором истечении срока действия ключей,
сгенерированных с квантово-уязвимыми параметрами.
Добавление временно"=зависимой оценки риска (с учётом горизонта
появления криптографически значимых квантовых компьютеров к 2030\,г.)
сделает приоритизацию плана миграции более реалистичной.
